%% 自動生成（gen.py）
\bktocchap{5-0}{スピードメーターは何を計算している？}{「今の速さ」は、どうやって決める？}{1}{2}{2}{8}{J}
\bktocpart{第1部}{微分}{16}
\bktocchap{5-1}{平均の速さと瞬間の速さ}{2時間で100 km走った車は、ずっと時速50 km？}{1}{2}{2}{18}{C}
\bktocchap{5-2}{時間をどんどん短くしてみよう}{0.1秒、0.01秒…と縮めると、ある値に近づく}{1}{2}{2}{40}{C}
\bktocchap{5-3}{かぎりなく近づく：極限と接線}{曲線にぴったりふれる直線の傾き}{1}{3}{3}{60}{C}
\bktocchap{5-4}{表から導関数を見つけよう}{表をながめると、微分のきまりが見えてくる}{1}{3}{3}{78}{C}
\bktocchap{5-5}{位置→速度→加速度}{2回微分すると、何がわかる？}{1}{3}{3}{104}{C}
\bktocchap{5-6}{微分の道具箱：積と合成関数}{歯車をつなぐと、回る速さはかけ算になる}{1}{4}{4}{126}{C}
\bktocchap{5-7}{いちばん高いのはどこ？　最大と最小}{いちばん高い点では、速度が0になる}{1}{3}{3}{154}{C}
\bktocchap{5-8}{短い時間なら、かけ算で近似できる}{「次の位置」を少しずつ予想する}{1}{2}{2}{174}{C}
\bktocpart{第2部}{積分}{194}
\bktocchap{5-9}{グラフの面積は進んだ距離}{速さのグラフの面積が、進んだ道のりになる}{2}{2}{2}{196}{C}
\bktocchap{5-10}{細い長方形に分けてたそう}{曲がった形の面積を、細い長方形で求める}{2}{2}{4}{228}{C}
\bktocchap{5-11}{積分は微分の逆}{面積がふえる速さは、その場所の高さ}{2}{3}{3}{262}{C}
\bktocchap{5-12}{物理での積分：仕事と流量}{ばねをのばす仕事は、三角形の面積}{2}{3}{3}{284}{C}
\bktocchap{5-13}{積分定数は「はじめの状態」}{スタート地点がわからないと、位置は決まらない}{2}{3}{3}{310}{C}
\bktocpart{第3部}{特別な関数}{332}
\bktocchap{5-14}{ふしぎな数 e}{微分しても形が変わらない、特別な関数}{3}{3}{4}{334}{C}
\bktocchap{5-15}{sin・cosを微分すると？}{微分すると、波が1/4周期ずれる}{4}{4}{4}{362}{C}
\bktocchap{5-16}{積分の道具箱：置換積分と部分積分}{むずかしい積分を、知っている形に変える}{2}{4}{4}{392}{C}
\bktocchap{5-17}{近くでは多項式で近似できる：テイラー展開}{小さくゆれる振り子は、$\sin\theta\approx\theta$で計算できる}{1}{4}{4}{420}{C}
\bktocchap{5-18}{まとめ：微分したときのふるまい}{微分方程式の「答えの候補」がそろった}{1}{3}{3}{444}{C}
\par\vspace{2mm}
\bktocbacklist{\bkbi{こたえ}{464}\bkbi{\R{記号}{き|ごう}カード}{477}\bkbi{よく\R{使}{つか}う\R{定数}{てい|すう}・\R{単位}{たん|い}・\R{公式}{こう|しき}}{480}\bkbi{できるようになったことチェック\R{表}{ひょう}}{484}\bkbi{つぎの\R{巻}{かん}へ}{487}\bkbi{さくいん}{489}}
